\chapter{Ранг-один угловое чтение семейной кривизны}

Предыдущий гейт оставил во второй степени только опорную линию
\(\mathbb C\rho\) и истолковал исчезновение \(Q\)-блока как недостаток для
полного отображения момента. Перекрёстный аудит проекта показывает, что это
требование было слишком сильным: полная семейная кривизна уже выведена в
модулярном чтении, тогда как ранг-один угол должен описывать наблюдение в
выбранном состоянии, а не заново хранить весь семейный оператор.

\section{Антикруговой журнал}

Этот гейт не повторяет:
\begin{enumerate}
 \item обычный спектральный квадрат, который не различает ориентацию;
 \item отображающий конус, сохраняющий крайний составной путь;
 \item фактор форм Конна, уже вычисленный в предыдущей главе;
 \item стандартную суперсвязность, снова дающую сумму матриц Грама.
\end{enumerate}
Используется более ранний результат проекта о производных секторных мерах:
один родительский след допускает петлевое суммирование, нормированное
угловое среднее и чистое состояние как разные заранее определённые типы
наблюдения.

\section{Полная и угловая кривизны}

После ограничения на \(S_4/A_4\)-эквивариантные стрелки
\begin{equation}
 X=rV,\qquad Y=\Phi I_3,\qquad VV^*=I_3
\end{equation}
средний блок модулярной кривизны равен
\begin{equation}
 \mathcal F_{\mathrm{fam}}
 =XX^*-Y^*Y
 =(r^2-|\Phi|^2)I_3
 =\mu I_3.
 \label{eq:v5-full-equivariant-moment}
\end{equation}
Пусть \(\rho\) --- уже выбранный проектор ранга один. Компрессия в его угол
имеет вид
\begin{equation}
 C_\rho(A)=\rho A\rho.
 \label{eq:v5-state-corner-compression}
\end{equation}
Она не унитальна как отображение в \(M_3\), поскольку \(C_\rho(I)=\rho\),
но унитальна как отображение в угловую алгебру \(\rho M_3\rho\) с единицей
\(\rho\). Карта вполне положительна, идемпотентна, совместима с инволюцией
и ковариантна при одновременном вращении \(A\) и \(\rho\).

Для кривизны~\eqref{eq:v5-full-equivariant-moment}
\begin{equation}
 C_\rho(\mu I_3)=\mu\rho.
 \label{eq:v5-corner-moment-readout}
\end{equation}
Это в точности направление \(\mathbb C\rho\), которое пережило
факторизацию двухформ в предыдущей главе.

\section{Почему коэффициент не потерян}

Полное семейное чтение использует нормированный след
\begin{equation}
 \tau_3((\mu I_3)^2)=\mu^2.
\end{equation}
Угловое чтение состояния использует условно нормированный след
\begin{equation}
 \tau_\rho(B)=\frac{\Tr(\rho B\rho)}{\Tr\rho},
 \qquad
 \tau_\rho((\mu\rho)^2)=\mu^2.
 \label{eq:v5-corner-normalized-norm}
\end{equation}
Следовательно, квадрат кривизны сохраняется без вставки множителя три.

На физическом пространстве
\(\mathbb C^3_{\mathrm{fam}}\otimes\mathbb C^{15}_{\mathrm{obs}}\)
проектор \(P_\rho=\rho\otimes I_{15}\) имеет ранг пятнадцать из сорока
пяти. Полный петлевой след поэтому даёт фиксированный множитель
\begin{equation}
 \tau_{45}(P_\rho\mu^2)=\frac13\mu^2,
\end{equation}
тогда как нормированный след угла снова даёт \(\mu^2\). Это не смена
нормировки после вычисления: первый функционал считает кратность состояний
в петле, второй относится к одному нормированному семейному состоянию.
Такое различие уже было заранее зафиксировано в гейтах производных мер и
единого следа \(M_{60}\).

\section{Совместимость с наблюдаемым чтением}

Компрессия \(C_\rho\) действует только на семейном множителе. Поэтому она
коммутирует с проекцией наблюдаемого множителя на точные блоки
\begin{equation}
 6+2+3+3+1=15
\end{equation}
Стандартной модели. Машинная проверка на общей эрмитовой матрице размера
сорок пять даёт нулевые, с точностью до округления, дефекты:
\begin{enumerate}
 \item совместимости с семейным частичным следом;
 \item совместимости с пятиблочным наблюдаемым ожиданием;
 \item совместимости условного наблюдаемого чтения состояния с тем же
 пятиблочным ожиданием.
\end{enumerate}

\begin{theorem}[Угловое чтение, а не повторное происхождение]
Одномерный фактор \(\mathbb C\rho\) является каноническим ранг-один чтением
уже выведенной эквивариантной семейной кривизны \(\mu I_3\). Нормированный
след угла сохраняет \(\mu^2\), полный родительский след сохраняет
фиксированную кратность \(1/3\), а точное чтение Стандартной модели
коммутирует с этой редукцией. Свободный относительный коэффициент не
возникает.
\end{theorem}

\begin{remark}
Гейт не утверждает, что полная матричная кривизна восстанавливается из
одного состояния. Редукция необратима и должна такой оставаться. Получен
согласованный переход между двумя типами наблюдения одного источника.
\end{remark}

\begin{protocol}
\begin{enumerate}
 \item полная модулярная кривизна: \textbf{взята из ранее доказанного гейта};
 \item её ранг-один компрессия: \textbf{получена};
 \item полная положительность угловой карты: \textbf{проверена};
 \item сохранение нормированного квадрата кривизны: \textbf{получено};
 \item множитель петлевой кратности \(1/3\): \textbf{выведен из ранга};
 \item совместимость с точным чтением Стандартной модели:
 \textbf{получена};
 \item новый вес или новая калибровочная группа: \textbf{не введены};
 \item динамика внедиагонального юкавского угла: \textbf{открыта};
 \item физическое замыкание: \textbf{не достигнуто}.
\end{enumerate}
\end{protocol}

Следующий гейт должен строить двустороннюю связность только для динамики
внедиагонального модуля. Перенос скалярной семейной кривизны между чтениями
уже замкнут угловой компрессией и не должен повторно усложняться.

Машинный сертификат сохранён в
\path{s2t_v5_state_corner_curvature_readout_gate.py} и
\path{s2t_v5_state_corner_curvature_readout_gate_results.json}.